\section{Additional Details}

\subsection{Framework Overview}

The end-to-end workflow is: real decode run $\rightarrow$ trace adapter $\rightarrow$ region-level trace $\rightarrow$ virtual PIM replay $\rightarrow$ policy comparison. The framework deliberately separates runtime measurement from the virtual backend so that different cost-model settings can be replayed on the same collected traces.

\subsection{Additional Tables}

Table~\ref{tab:main-modern-results} summarizes representative modern-family results, and Table~\ref{tab:family-onset} summarizes the first positive context by family.

\begin{table}[t]
\centering
\caption{Representative regime-transition results on modern 4-bit instruct models.}
\label{tab:main-modern-results}
\begin{tabular}{lcccc}
\toprule
Model & Context & Online & KVRegime & Oracle \\
\midrule
Qwen2.5-1.5B & 256 & 1.000 & 1.000 & 1.274 \\
Qwen2.5-1.5B & 512 & 1.169 & 1.169 & 1.306 \\
Qwen3-8B & 256 & 1.155 & 1.155 & 1.277 \\
SmolLM2-1.7B & 256 & 1.150 & 1.150 & 1.276 \\
Phi-3.5-mini & 256 & 1.140 & 1.140 & 1.276 \\
Phi-3.5-mini & 384 & 1.181 & 1.181 & 1.297 \\
\bottomrule
\end{tabular}
\end{table}

\begin{table}[t]
\centering
\caption{Family-dependent onset of the positive offload regime.}
\label{tab:family-onset}
\begin{tabular}{lc}
\toprule
Family & First online-positive context \\
\midrule
gpt2 & 512 \\
opt-125m & 512 \\
Qwen2.5-1.5B & 384 \\
Qwen3-8B & 256 \\
SmolLM2-1.7B & 256 \\
Phi-3.5-mini & 256 \\
\bottomrule
\end{tabular}
\end{table}


\subsection{Robustness Notes}

The sensitivity study changes PIM bandwidth and latency presets while leaving the decode traces fixed. The most robust conclusion is the long-context positive-regime story. Borderline early-onset short-context points are more parameter-sensitive, which is why the main paper phrases them as family-dependent onset evidence rather than as universal short-context gains.

\subsection{Reproducibility}

All main figures in the paper are generated from reproducible scripts and stored under the project \texttt{final/} package. The trace schema, simulator, cost model, and policy implementations are contained in the repository. The paper focuses on runtime-grounded virtual replay rather than private serving traces or proprietary hardware measurements.
